\section{Kết luận}\label{sec:concl}
Báo cáo trả lời bốn câu hỏi nghiên cứu như sau. \textbf{RQ1:} trên cùng split và cùng 100 ảnh ứng viên, QA-Graph đạt MRR $64{,}89$ trên VG$\cap$COCO và $64{,}41$ trên PSG$\cap$COCO (toàn bộ ảnh test), so với $60{,}40$ và $59{,}59$ của CLIP giữ nguyên trọng số, và hơn mốc CLIP + adapter khoảng 1,4--1,9 điểm. \textbf{RQ2:} xáo cạnh kéo kết quả xuống dưới mốc adapter, nên cải thiện đến từ nội dung cạnh chứ không từ tham số học thêm; nhưng gần như toàn bộ đến từ cạnh \texttt{contains}, còn quan hệ scene graph và ConceptNet không thêm đáng kể. Trên VG mức tăng tập trung ở truy vấn trực tiếp, và cạnh quan hệ chỉ chuyển một phần nhỏ lợi ích sang nhóm gián tiếp mà không tăng tổng; nhóm gián tiếp của PSG có tăng nhưng cũng chỉ nhờ \texttt{contains}. \textbf{RQ3:} AdaMF-MAT tái lập được trên MKG-W và MKG-Y (sai khác dưới 1 điểm), nhưng khi áp dụng vào truy vấn ảnh không vượt baseline khớp tên đối tượng trên cả hai bộ dữ liệu. \textbf{RQ4:} trên 34 truy vấn trực tiếp đã duyệt, xóa cạnh giải thích làm điểm ảnh đứng đầu giảm nhiều hơn xóa cạnh đối chứng, nên giải thích trung thực với hàm chấm điểm; tuy vậy các cạnh được chọn chủ yếu nối nhãn với ảnh, và độ đúng thị giác của từng cạnh cũng như khả năng giải thích quan hệ nhiều bước cần được kiểm chứng tiếp.
